% !PRINTED 758
\weberpage{758}

If we now set $x=0,\alpha_1,\alpha_2,\ldots,\alpha_\mu$ in equation (12) of the preceding section and apply equation (3), we obtain:
\weq{4}{0 = CF(0) + F(\alpha_1) + F(\alpha_2) + \cdots + F(\alpha_\mu) + CQ(0) + e^{a_1}Q(\alpha_1) + e^{a_2}Q(\alpha_2) + \cdots + e^{a_\mu}Q(\alpha_\mu),}
and we have to demonstrate the impossibility of such an equation for a suitable choice of $f(x)$.

Let $p$ again be a prime determined to increase without limit, and
\weq{5}{f(x) = \dfrac{a^{\mu p+p-1}x^{p-1}(x-\alpha_1)^p(x-\alpha_2)^p\ldots(x-\alpha_\mu)^p}{\Pi(p-1)},}
so that $f(x)$ is a function with rational coefficients.

We now form, as in the preceding section, by arranging according to powers of $x$:
\[
f(x) = \dfrac{A_{p-1}x^{p-1} + A_px^p + A_{p+1}x^{p+1}\ldots}{\Pi(p-1)},
\]
in which the $A_{p-1},A_p,\ldots$ are integral rational numbers, and
\[
A_{p-1} = (-1)^\mu a^{\mu p+p-1}\alpha_1^p\alpha_2^p\ldots\alpha_\mu^p.
\]
If we therefore assume $p$ to be greater than each of the integers
\[
a, a^\mu\alpha_1\alpha_2\ldots\alpha_\mu
\]
then $A_{p-1}$ is not divisible by $p$, and we have
\[
F(0) = A_{p-1} + pA_p + p(p+1)A_{p+1} + \cdots,
\]
i.\,e. an integer not divisible by $p$. Likewise, if $p$ is taken large enough, $C$ is not divisible by $p$.

On the other hand, we arrange the numerator of $f(x)$ according to powers of $a(x-\alpha_1)$ and obtain
\[
f(x) = \dfrac{B_pa^p(x-\alpha_1)^p + B_{p+1}a^{p+1}(x-\alpha_1)^{p+1} + \cdots}{\Pi(p-1)},
\]
in which the $B_p,B_{p+1},\ldots$ are indeed no longer rational, but are nevertheless \gesperrt{integral algebraic numbers}, since the numerator of $f(x)$ is a function with integral coefficients in $ax,a\alpha_1,\ldots,a\alpha_\mu$.

From this it now follows, as in the preceding section:
\[
F(\alpha_1) = pB_pa^p + p(p+1)B_{p+1}a^{p+1} + \cdots
\]
If one forms $F(\alpha_2),\ldots,F(\alpha_\mu)$ in the same way, and observes that the sum $B_p(\alpha_1)+B_p(\alpha_2)+\cdots+B_p(\alpha_\mu)$ is an integral rational number, it follows that
\[
F(\alpha_1) + F(\alpha_2) + \cdots + F(\alpha_\mu)
\]
%CONT>>
